% ============ 第 7 讲 赫克歇尔-俄林模型 ============
\chap{第 7 讲　赫克歇尔--俄林模型}{The Heckscher--Ohlin Model}

\textbf{本讲导读}：本讲对应 Lecture 7（Heckscher--Ohlin Model, Oct.\ 14），共 69 页（含附录与 3 页课后手写笔记）。内容：H-O 模型的六条假设（要素自由流动、部门要素密集度差异、国家要素丰裕度差异、自由贸易、技术相同、同质位似偏好）；三个核心概念（国家要素丰裕度、跨行业要素密集度、行业内的要素比例）；相对要素需求曲线的推导；弯曲 PPF 的形状与偏斜；自给自足均衡、贸易模式与 H-O 定理；贸易的要素价格效应（Stolper--Samuelson 定理）；要素价格均等化；附录 1（PPF 斜率）与附录 2（方程结构）；课后手写笔记给出 SFM 与 H-O 要素市场的系统对照。

\secn{7.1 从特定要素模型到 H-O：动机与保留}

\textbf{特定要素模型要点与局限（第 2--3 页）}：贸易模式由\kw{比较优势}决定；两国总体都\kw{得益}；福利分配——\redkw{Gain} 给被困在比较优势部门的要素所有者、\redkw{Lose} 给被困在比较劣势部门的要素所有者、\redkw{Ambiguous} 给可流动要素（取决于消费模式）。局限（第 3 页，高亮）：
\begin{itemize}
  \item \kw{比较优势的决定因素似乎很复杂}（红底高亮）——板书：$\frac{P_M}{P_A} = \frac{a_A}{a_M}\left( \frac{T}{K} \cdot \frac{L_M}{L_A} \right)^{1/3}$，其中 $\frac{L_M}{L_A}$ 是\kw{内生配置}（手写 Factor Endowment / Productivity / Endogenous Allocation 标注）；
  \item \kw{要素可流动时福利含义模糊}（黄底高亮）；
  \item 福利影响可以超出特定部门（如美国的非熟练工人）；
  \item \kw{短期 vs.\ 长期}（粉底高亮）：短期被困、长期可以调整。
\end{itemize}

\textbf{H-O 模型预览（第 4 页）}：从 SFM 保留的关键特征——\redkw{比较优势}决定专业化；多要素带来的弯曲 PPF；福利含义：不同要素的所有者。\kw{新设定}：\term{要素跨部门完全流动}；但\term{以不同比例被使用}（used in different proportions）。

\secn{7.2 模型设定：六条假设与三个核心概念}

\begin{assumpbox}[Heckscher--Ohlin 模型设定（第 6、11--12、23 页）]
\begin{itemize}
  \item \term{Assumption 1}：两种生产要素——\kw{资本（K）}与\kw{劳动（L）}（粉底高亮），可以在行业间\term{自由流动}；
  \item \term{Assumption 2}：两个部门——\kw{鞋（S）}与\kw{计算机（C）}（粉底高亮），鞋的生产是``\redkw{劳动密集（labor intensive）}''（高亮）；
  \item \term{Assumption 3}：Foreign 是``\redkw{劳动丰裕}''，Home 是``\redkw{资本丰裕}''（红字+高亮）——定义：Foreign 劳动丰裕指 Foreign 的劳动--资本\term{禀赋比}更高：
  \[
  {\color{hlorange}\frac{\bar{L}^*}{\bar{K}^*} > \frac{\bar{L}}{\bar{K}}}
  \]
  （注意是\term{相对}禀赋，不是绝对量！）；
  \item \term{Assumption 4}：商品可以自由贸易，但\term{劳动与资本不能跨国流动}；
  \item \term{Assumption 5}：两国生产两种商品的\term{技术相同}；
  \item \term{Assumption 6}：两国消费者\term{偏好相同}，且对计算机 vs.\ 鞋的相对需求不依赖于收入水平——``\kw{identical \& homothetic}''（高亮；手写批注：[Assumption 6] ``Homothetic Preferences'' $\Rightarrow \frac{MU_1}{MU_2} \sim \frac{C_1}{C_2}$——收入扩张线为过原点直线）。
\end{itemize}
\end{assumpbox}

\begin{defbox}[三个核心概念（第 7、13 页，True or False 练习）]
\begin{enumerate}
  \item \kw{国家要素丰裕度（Country factor abundance）}（高亮）：``Home 资本丰裕 $=$ Home 拥有比 Foreign 更多的资本？''——\redkw{False}！手写：``relatively''、$\frac{\bar{L}}{\bar{K}}$——是\term{相对禀赋比}的比较；
  \item \kw{跨行业要素密集度（Factor intensity across industries）}（高亮）：``计算机资本密集 $=$ 计算机使用的资本总量比鞋多？''——\redkw{False}！手写：``relative demand curve''——是在\term{同一} $w/r$ 下 $K/L$ 的比较（相对要素需求曲线的位置）；
  \item \kw{行业内要素比例（Factor proportions within an industry）}（高亮）：``更高的 $w/r$ 使每个行业的厂商选择更高的 $K/L$''——\redkw{True}！手写：$\frac{w}{r} \uparrow \Rightarrow \frac{L}{K} \downarrow$；``curve slope''（沿相对要素需求曲线移动）。
\end{enumerate}
\end{defbox}

\textbf{要素密集度的定义与决定（第 7 页）}：鞋的生产是``\redkw{劳动密集}''（红字下划线）指：生产鞋比生产计算机每单位资本使用更多劳动，即 \kw{$L_S/K_S > L_C/K_C$}（粉底高亮）——\redkw{under what condition?}（红字设问：在相同的 $w/r$ 下！）什么决定部门内 K 与 L 的使用？\kw{生产技术}（高亮）；\kw{要素的相对价格}（高亮）：$w/r$ 更高 $\Rightarrow$ 相对 K 而言更少地使用 L。\term{Wage（w）}：支付给劳动；\term{Rental rate（r）}：支付给 K（红字）。

\textbf{相对要素需求曲线（第 8--9 页）}：\kw{相对要素需求（Relative factor demand，高亮）}：假设鞋的劳动/资本比在任何给定的 $w/r$ 下都高于计算机：$L_S/K_S > L_C/K_C$ for any given $w/r$（蓝字）——\kw{Assumption 2 意味着两条曲线永不相交}（蓝字）。图~\ref{fig:ch07-rfd}：横轴为各行业的劳动/资本比、纵轴为 wage/rental；手写批注：``\kw{Factor intensity}: $\frac{L_S}{K_S}$ curve is on the right''（鞋的曲线在右边）。

\begin{figure}[H]
  \centering
  \includegraphics[width=0.55\linewidth]{figs/ch07_fig_01.png}
  \caption{第 8 页：两行业的相对要素需求曲线（鞋 $L_S/K_S$ 在计算机 $L_C/K_C$ 右侧，永不相交）；手写批注：``Factor intensity: $\frac{L_S}{K_S}$ curve is on the right''；沿曲线的移动由 $w/r$ 变化驱动。}
  \label{fig:ch07-rfd}
\end{figure}

\textbf{向下倾斜的推导（第 9 页）}：由最优生产决策条件（FOCs）导出：板书 $w = P_S \cdot MPL_S$ 且 $r = P_S \cdot MPK_S$（高亮）$\Rightarrow$
\[
{\color{hlorange}\frac{w}{r} = \frac{MPL_S}{MPK_S}}
\]
（红框；手写箭头连接）。$MPL_S/MPK_S$ 取决于 $L_S/K_S$；更多 L、更少 K：更高的 MPK 与更低的 MPL（报酬递减）。例如 C--D 生产函数。手写：$\frac{L}{K} \uparrow \Rightarrow MPK \uparrow$；$\frac{K}{L} \downarrow \Rightarrow MPL \downarrow$。

\textbf{数值例子：要素密集度（第 10 页，图~\ref{fig:ch07-intensity}）}：从 $w/r = MPL/MPK$ 推导 $L_S/K_S$ 与 $L_C/K_C$。生产函数 $Y_S = a_S K_S^{\beta_S} L_S^{1-\beta_S}$，$Y_C = a_C K_C^{\beta_C} L_C^{1-\beta_C}$，且 $1 > \beta_C > \beta_S > 0$（$\beta$ 是资本份额——$\beta_C$ 大 $=$ 计算机更资本密集）。板书推导：
\[
\frac{w}{r} = \frac{MPL_g}{MPK_g} = \frac{(1-\beta_g) a_g K_g^{\beta_g} L_g^{-\beta_g}}{\beta_g a_g K_g^{\beta_g - 1} L_g^{1-\beta_g}} = \frac{1-\beta_g}{\beta_g} \cdot \left( \frac{L_g}{K_g} \right)^{-1}
\;\Longrightarrow\;
{\color{hlorange}L_g / K_g = \frac{1-\beta_g}{\beta_g} \left( \frac{w}{r} \right)^{-1}}
\]
（手写注：$\beta$ 越大、$\frac{1}{\beta_g} - 1$ 越小）。结论：$L_S/K_S$ 与 $L_C/K_C$ 都随 $w/r$ 递减（沿曲线移动）；且对一切 $w/r$ 有 $L_S/K_S > L_C/K_C$——手写：``S is labor intensive: $\beta_S < \beta_C$''。

\begin{figure}[H]
  \centering
  \includegraphics[width=0.66\linewidth]{figs/ch07_fig_02.png}
  \caption{第 10 页板书：由 $w/r = MPL_g/MPK_g$ 推导 $L_g/K_g = \frac{1-\beta_g}{\beta_g} (\frac{w}{r})^{-1}$；$L_S/K_S$、$L_C/K_C$ 都随 $w/r$ 递减，且 $L_S/K_S > L_C/K_C$ 对一切 $w/r$ 成立（``S is labor intensive: $\beta_S < \beta_C$''）。}
  \label{fig:ch07-intensity}
\end{figure}

\secn{7.3 PPF 的形状：弯曲与偏斜}

\textbf{从两组方程解出 PPF（第 14 页）}：生产函数 $Y_S = a_S K_S^{\beta_S} L_S^{1-\beta_S}$、$Y_C = a_C K_C^{\beta_C} L_C^{1-\beta_C}$；资源约束 {\color{hlorange}$L_S + L_C = \bar{L}$}、{\color{hlorange}$K_S + K_C = \bar{K}$}（高亮）。Foreign 类似。提醒：Home 资本丰裕，$\bar{K}/\bar{L} > \bar{K}^*/\bar{L}^*$。两国 PPF 形状如何？

\textbf{PPF 斜率（第 15 页，图~\ref{fig:ch07-ppfslope}）}：PPF 绝对斜率 $= P_C/P_S$（贸易下一价定律）。由两个要素市场均衡（均高亮）：
\[
\text{劳动市场：}\; P_S MPL_S = w = P_C MPL_C \;\Longrightarrow\; \frac{P_C}{P_S} = \frac{MPL_S}{MPL_C},
\]
\[
\text{资本市场：}\; P_S MPK_S = r = P_C MPK_C \;\Longrightarrow\; \frac{P_C}{P_S} = \frac{MPK_S}{MPK_C}.
\]
板书（手写）：$\frac{MPK_S}{MPK_C} = \frac{P_C}{P_S} = \frac{MPL_S}{MPL_C} = $ Slope。手写注（增加 C 的产量）：$K_C \uparrow \Rightarrow MPK_C \downarrow, MPL_C \uparrow$；$L_C \uparrow \Rightarrow MPK_C \uparrow, MPL_C \downarrow$。

\begin{figure}[H]
  \centering
  \includegraphics[width=0.66\linewidth]{figs/ch07_fig_03.png}
  \caption{第 15 页板书：PPF 斜率 $= \frac{MPK_S}{MPK_C} = \frac{P_C}{P_S} = \frac{MPL_S}{MPL_C}$；两个要素市场均衡条件（$P_S MPL_S = w = P_C MPL_C$、$P_S MPK_S = r = P_C MPK_C$）。}
  \label{fig:ch07-ppfslope}
\end{figure}

\textbf{PPF 为何弯曲（第 16 页，图~\ref{fig:ch07-ppfbend}）}：\kw{PPF 斜率在横轴产量增加时变陡}，除非两行业要素密集度相同。直觉（蓝字+板书）：向右（向下）移动时——开始时资本在鞋业相对充裕，把资源转入计算机起初只牺牲相对少的鞋（Point A，板书标注）；随着计算机生产扩张，鞋业能释放的资本越来越少，再生产一台计算机要牺牲更多鞋（Point B）。（正式推导见附录。）

\begin{figure}[H]
  \centering
  \includegraphics[width=0.4\linewidth]{figs/ch07_fig_04.png}
  \caption{第 16 页板书：沿 PPF 从 S 端向 C 端移动——$A$ 点（$\Delta S$ 小、$\Delta C$ 大）到 $B$ 点（$\Delta S$ 大、$\Delta C$ 小）：机会成本递增。}
  \label{fig:ch07-ppfbend}
\end{figure}

\textbf{PPF 的偏斜（第 17 页）}：Home 资本丰裕（$\bar{K}/\bar{L} > \bar{K}^*/\bar{L}^*$，高亮；手写：``$\Rightarrow$ PPF 的高低/胖瘦''）$\Rightarrow$ \term{Home 的 PPF 更偏向计算机}（资本密集品）。图~\ref{fig:ch07-skew}：(a) Home（红，偏 C 轴）；(b) Foreign（洋红，偏 S 轴）。

\begin{figure}[H]
  \centering
  \includegraphics[width=0.72\linewidth]{figs/ch07_fig_05.png}
  \caption{第 17 页：两国 PPF 的偏斜——资本丰裕的 Home（左）偏向计算机（C）；劳动丰裕的 Foreign（右）偏向鞋（S）。}
  \label{fig:ch07-skew}
\end{figure}

\secn{7.4 自给自足均衡、贸易模式与 H-O 定理}

\textbf{自给自足均衡（第 19--20 页，图~\ref{fig:ch07-autarky}）}：无差异曲线与 PPF 相切；Home 的计算机相对价格比 Foreign \kw{更低}（板书：Home 图 $C$ 大、$S$ 小，$\frac{P_C}{P_S} \downarrow$；Foreign 图 $C$ 小、$S$ 大，$\frac{P_C^*}{P_S^*}$ 大）——手写批注顶部：[Assumption 6] ``Homothetic Preferences'' $\Rightarrow \frac{MU_1}{MU_2} \sim \frac{C_1}{C_2}$。\kw{Home 在计算机上有比较优势}（高亮）。

\begin{figure}[H]
  \centering
  \includegraphics[width=0.7\linewidth]{figs/ch07_fig_06.png}
  \caption{第 19--20 页：两国自给自足均衡——Home（左）生产更多 C、价格线更平（$\frac{P_C}{P_S}$ 小）；Foreign（右）相反。}
  \label{fig:ch07-autarky}
\end{figure}

\textbf{贸易模式（第 21--22 页）}：自给自足下 $\left( \frac{P_C}{P_S} \right)^A < \left( \frac{P_C^*}{P_S^*} \right)^A$；Home 有 CA in 计算机。贸易使相对价格均等：
\[
\left( \frac{P_C}{P_S} \right)^A < \left( \frac{P_C}{P_S} \right)^W < \left( \frac{P_C^*}{P_S^*} \right)^A .
\]
\redkw{贸易条件}（出口品的相对价格）对两国都上升——与李嘉图模型一样。Home 出口计算机——\kw{密集使用}（下划线）其\term{相对丰裕}要素（K）的商品；Foreign 出口鞋——密集使用其相对丰裕要素（L）的商品。这就是\redkw{Heckscher--Ohlin Theorem}（红字）。正式证明：见下一讲。

\begin{thmbox}[H-O 定理（第 22--23 页）]
\kw{一国出口密集使用其相对丰裕要素生产的商品。}所需假设回顾（第 23 页）：A1：L 与 K 跨行业流动；A2：鞋的生产相对计算机是劳动\term{密集}的（K intensive for C）；A3：两国 L 与 K 的数量不同，Foreign \term{丰裕}于 L、Home 丰裕于 K；A4：商品自由贸易；A5\&6：两国的生产技术与消费偏好相同。
\end{thmbox}

\textbf{复习填空（第 24 页板书）}：Home is \kw{K} abundant. Computer production is \kw{K} intensive. Home's autarky $\frac{P_C}{P_S}$ is \kw{$<$} than Foreign's. With trade, $\frac{P_C}{P_S}$ \kw{$>$} at Home. Home produces more \kw{computer} than it consumes and exports \kw{computer}（板书填空）。

\textbf{贸易得益（第 25 页，图~\ref{fig:ch07-gains}）}：像 SFM 一样\term{总体上有贸易得益}；要素所有者之间的分配效应如何？图：世界相对价格高于 Home 自给自足（红字 World rel.\ price $>$ Home autarky），高于 Foreign 自给自足的价格对其反向；手写：``gains from trade''。

\begin{figure}[H]
  \centering
  \includegraphics[width=0.72\linewidth]{figs/ch07_fig_07.png}
  \caption{第 25 页：贸易得益——世界相对价格线（红）比两国自给自足价格线都更陡/更平（各自朝出口品方向），消费点移到更高的无差异曲线上。}
  \label{fig:ch07-gains}
\end{figure}

\secn{7.5 贸易对要素价格的影响：Stolper--Samuelson 定理}

\textbf{分析框架（第 27--28 页）}：在 H-O 框架内：贸易对\kw{真实工资}（$w/P_S$、$w/P_C$——手写注 $MPL_S$、$MPL_C$）与\kw{真实租金率}（$r/P_S$、$r/P_C$——手写注 $MPK_S$、$MPK_C$）的影响？先看\term{相对要素价格}（$w/r$）再看真实值。设 Home 有 \kw{12 单位资本与 20 名工人}（板书蓝底高亮），初始 $\frac{w}{r} = 1$（Autarky，手写）：

\begin{figure}[H]
  \centering
  \includegraphics[width=0.72\linewidth]{figs/ch07_fig_08.png}
  \caption{第 28 页：要素再配置的数值例子——初始配置：计算机 $K_C = 6$、$L_C = 2$（$L/K = 1/3$）；鞋 $K_S = 6$、$L_S = 18$（$L/K = 3$）。贸易使计算机扩张、鞋收缩（板书：$K_C: 6 \to 9$，$K_S: 6 \to 3$）。设问：要素比例能否保持不变？}
  \label{fig:ch07-realloc}
\end{figure}

\textbf{旧要素比例下劳动过剩（第 29 页，图~\ref{fig:ch07-old}）}：保持 $\frac{w}{r} = 1$ 与各行业原有要素比例（$\frac{L}{K}$：C 中 $1/3$、S 中 3，手写）：计算机需要 $9 \times (1/3) = 3$ 名工人、鞋需要 $3 \times 3 = 9$ 名；合计只需 12 名，\redkw{旧要素比例使 8 名工人失业}（板书手写箭头）。\kw{工资相对租金率必须发生什么变化？}（蓝字设问）

\begin{figure}[H]
  \centering
  \includegraphics[width=0.72\linewidth]{figs/ch07_fig_09.png}
  \caption{第 29 页：旧要素比例下只需要 12 名工人（20 名可用）——8 名工人闲置，$w/r$ 必须下降。}
  \label{fig:ch07-old}
\end{figure}

\textbf{要素比例的调整（第 30 页，图~\ref{fig:ch07-adjust}）}：为吸收多余劳动，两个行业都需要更高的 $L/K$；使厂商愿意这样做，\kw{$w/r$ 必须下降}。新配置（板书）：$\frac{w}{r} = 0.6$；计算机 $K = 9$、$L/K = 5/9$、$L = 5$；鞋 $K = 3$、$L/K = 5$、$L = 15$——全部 12 单位资本与 20 名工人就业。\kw{为什么两个行业比例都上升而国家 L/K 不变？}（蓝字设问——因为资本从劳动密集的鞋业流向资本密集的计算机业，两行业的 $L/K$ 都向国家平均靠拢的方向移动。）

\begin{figure}[H]
  \centering
  \includegraphics[width=0.72\linewidth]{figs/ch07_fig_10.png}
  \caption{第 30 页：新均衡配置——$w/r$ 从 1 降到 0.6，两行业 $L/K$ 都上升（$5/9$ 与 $5$），全部要素就业。}
  \label{fig:ch07-adjust}
\end{figure}

\textbf{相对要素供给与需求的正式框架（第 31--34 页）}：把数值例子形式化为相对供求框架。Home 劳动的\kw{相对供给}（supply，红底高亮）：
\[
\frac{\bar{L}}{\bar{K}} = \frac{L_S + L_C}{K_S + K_C}.
\]
\kw{相对需求}（demand，绿底高亮）：替换 $K_S + K_C = \bar{K}$（手写分解 $\frac{\bar L}{\bar K} = \frac{L_S}{\bar K} + \frac{L_C}{\bar K}$，利用 $\frac{K_S}{K_S} = 1$、$\frac{K_C}{K_C} = 1$，板书两步）：
\[
\frac{\bar{L}}{\bar{K}} = \frac{L_C}{K_C} \cdot \frac{K_C}{\bar{K}} + \frac{L_S}{K_S} \cdot \frac{K_S}{\bar{K}}
\]
（手写注：``Relative factor demand of each sector（pp.\ 7--10）'' $\to$ ``\kw{Weights of each sector}''）。LHS 是\term{相对供给}，由要素禀赋决定（垂直线）；RHS 是\term{相对需求}，由以下决定：劳动密集度（$w/r$ 的函数）：$L_C/K_C$ 与 $L_S/K_S$（浅蓝曲线）；各行业在资本使用中的份额：$K_C/\bar{K}$ 与 $K_S/\bar{K}$（$= 1 - K_C/\bar{K}$，权重）。手写板书（第 34 页）：Relative supply $=$ Fixed vertical line；Relative demand $=$ avg.\ labor intensities, weighted by capital share——\term{劳动的相对需求 $=$ 以各行业资本份额加权的劳动密集度平均}。图~\ref{fig:ch07-rsrd}：$RD$ 位于 $L_C/K_C$ 与 $L_S/K_S$ 之间（红蓝粗线）。

\begin{figure}[H]
  \centering
  \includegraphics[width=0.72\linewidth]{figs/ch07_fig_11.png}
  \caption{第 34 页：经济总体的相对要素供求——相对供给为垂直线 $\bar{L}/\bar{K}$；相对需求 $RD$ 是两行业劳动密集度按资本份额的加权平均，位于 $L_C/K_C$ 与 $L_S/K_S$ 之间。}
  \label{fig:ch07-rsrd}
\end{figure}

\textbf{贸易对相对要素价格的影响（第 35--38 页，图~\ref{fig:ch07-stolper}）}：Home 计算机行业扩张（贸易开放使 $P_C/P_S \uparrow$：生产点 $A \to B$ 沿 PPF 移动；手写批注：``① 判断贸易方向 $\to$ 部门扩张/收缩''）。四步（板书红字 ①②③④）：
\begin{enumerate}
  \item 计算机相对价格上升使相对需求曲线从 $RD_1$ 移到 $RD_2$（权重变化，手写``\kw{权重变化、Relative demand 移动}''）——因为 $K_C/\bar{K} \uparrow$（+）、$K_S/\bar{K} \downarrow$（−）（板书圈注）；
  \item 相对工资从 $(W/R)_1$ 降到 $(W/R)_2$（手写：``③ 新的均衡、$\frac{w}{r}$ 下降''；旁注 Numeric Example 对应 A $\to$ B 点）；
  \item 在新的相对工资下，每个行业的劳动/资本比都上升（手写：``④ 两个部门 $\frac{L}{K} \uparrow$''）。
\end{enumerate}

\begin{figure}[H]
  \centering
  \includegraphics[width=0.7\linewidth]{figs/ch07_fig_12.png}
  \caption{第 36--38 页：贸易的相对要素价格效应三步图——$RD$ 左移 $\Rightarrow$ $w/r$ 下降 $\Rightarrow$ 两行业 $L/K$ 都上升（沿各自曲线移动）。}
  \label{fig:ch07-stolper}
\end{figure}

\textbf{解析验证（第 39 页，图~\ref{fig:ch07-numeric}）}：代入相对要素需求 $L_i/K_i = \left( \frac{1-\beta_i}{\beta_i} \right) \left( \frac{w}{r} \right)^{-1}$（$i = C, S$；粉底）：
\[
\frac{\bar{L}}{\bar{K}} = \left( \frac{1-\beta_C}{\beta_C} \right) \left( \frac{w}{r} \right)^{-1} \frac{K_C}{\bar{K}} + \left( \frac{1-\beta_S}{\beta_S} \right) \left( \frac{w}{r} \right)^{-1} \frac{K_S}{\bar{K}}
\;\Longrightarrow\;
\frac{w}{r} = \frac{1}{\bar{L}} \left[ \frac{1-\beta_C}{\beta_C} \cdot K_C + \frac{1-\beta_S}{\beta_S} \cdot K_S \right]
\]
（手写：$K_C +\Delta K$、$K_S -\Delta K$）。C 扩张、S 收缩时 $\Delta \left( \frac{w}{r} \right) = \frac{\Delta K}{\bar{L}} \left[ \frac{1-\beta_C}{\beta_C} - \frac{1-\beta_S}{\beta_S} \right] < 0$（手写 $(−)$ 与``下降''标注）——因为 $\beta_C > \beta_S \Rightarrow \frac{1-\beta_C}{\beta_C} < \frac{1-\beta_S}{\beta_S}$。同时 $L_C/K_C \uparrow$、$L_S/K_S \uparrow$（劳动密集度在两个部门都上升）。

\begin{figure}[H]
  \centering
  \includegraphics[width=0.66\linewidth]{figs/ch07_fig_13.png}
  \caption{第 39 页板书：解析推导——C 扩张 S 收缩 $\Rightarrow$ $\frac{w}{r}$ 下降（$\Delta(w/r) = \frac{\Delta K}{\bar L}[\frac{1-\beta_C}{\beta_C} - \frac{1-\beta_S}{\beta_S}] < 0$），两部门劳动密集度都上升。}
  \label{fig:ch07-numeric}
\end{figure}

\textbf{Home 开放贸易的五步链条（第 40 页）}：手写（C $\uparrow$、$\frac{w}{r} \downarrow$）：① 计算机相对价格上升；② 计算机行业扩张；③ 在旧的 $w/r$ 下劳动相对需求下降；④ \kw{工资/租金率之比下降}；⑤ 每个行业的劳动密集度上升。

\textbf{真实要素回报（第 41--45 页）}：
\begin{itemize}
  \item \term{$w/r$ 下降}削弱劳动相对资本的收益；国家总体得益，收入的更大份额流向资本——\kw{资本所有者明确得益}（高亮）；\term{工人呢？可能得益更少或受损}（手写下划线）。
  \item \term{对 $MPK$ 与 $MPL$ 的影响（第 42--44 页板书）}：$MPK_C$ \kw{$\uparrow$} 因为 $L_C/K_C$ 上升（C 行业有更多工人操作机器）；$MPK_S$ \kw{$\uparrow$} 因为 $L_S/K_S$ 上升；$MPL_C$ \kw{$\downarrow$} 因为 $K_C/L_C$ 下降（C 行业每个工人的机器变少）；$MPL_S$ \kw{$\downarrow$} 因为 $K_S/L_S$ 下降。
  \item \term{租金率（第 43 页）}：$r = P_C \cdot MPK_C = P_S \cdot MPK_S$；真实租金率 $r/P_C = MPK_C \uparrow$、$r/P_S = MPK_S \uparrow$——\kw{租金率比任何一种商品价格涨得都快，资本所有者变好}（高亮）。
  \item \term{工资（第 45 页）}：$w = P_C \cdot MPL_C = P_S \cdot MPL_S$；真实工资 $w/P_C = MPL_C \downarrow$（$K_C/L_C$ 下降）、$w/P_S = MPL_S \downarrow$（$K_S/L_S$ 下降）——\kw{工资比任何一种商品价格涨得都慢，劳动所有者变差}（高亮）。
\end{itemize}

\textbf{名义价格的比例变化不等式（第 46 页，图~\ref{ch07-ineq}）}：与 SFM 一样可以推出名义比例变化的``不等式串''（手写：$\frac{w}{P_g} \downarrow$、$\frac{r}{P_g} \uparrow$）：
\[
{\color{warnred}\frac{\Delta w}{w} < \frac{\Delta P_S}{P_S} < \frac{\Delta P_C}{P_C} < \frac{\Delta r}{r}}
\]
（$P_C/P_S \uparrow$，黄底）。Foreign（$P_C/P_S$ 下降）情形：
\[
\frac{\Delta r}{r} < \frac{\Delta P_C}{P_C} < \frac{\Delta P_S}{P_S} < \frac{\Delta w}{w}.
\]

\begin{figure}[H]
  \centering
  \includegraphics[width=0.66\linewidth]{figs/ch07_fig_14.png}
  \caption{第 46 页：名义比例变化的不等式串（Home 情形：$\frac{\Delta w}{w} < \frac{\Delta P_S}{P_S} < \frac{\Delta P_C}{P_C} < \frac{\Delta r}{r}$）——工资涨幅低于两种商品价格（真实工资下降），租金涨幅高于两种商品价格（真实租金上升）。}
  \label{ch07-ineq}
\end{figure}

\begin{thmbox}[Stolper--Samuelson 定理（第 47 页，标题荧光高亮）]
当两种商品都被生产且要素可流动时，\kw{一种商品相对价格的上升}：
\begin{itemize}
  \item \term{提高}密集用于该商品生产的要素的\term{真实收益}；
  \item \term{降低}另一种要素的真实收益。
\end{itemize}
（PS：无论该要素被哪个行业雇用——要素的身份而非所在行业决定其得与失！）
\end{thmbox}

\textbf{贸易的福利效应（第 48 页，手写逻辑链）}：板书：\kw{要素丰裕（K）} $\xrightarrow{\text{H-O 定理}}$ \kw{出口部门（C）} $\longrightarrow$ \kw{密集使用要素的真实回报 $\uparrow$}（手写：$K: \frac{r}{P_C} \cdot \frac{r}{P_S}$）。若 Home 与更劳动丰裕的国家贸易：资本密集品相对价格上升；租金率比工资涨得快；\kw{资本所有者从贸易中得益}；\kw{工人从贸易中受损}；\kw{总体贸易得益：赢家可以补偿输家}（高亮）。

\secn{7.6 要素价格均等化（Factor Price Equalization）}

\textbf{回顾禀赋（第 50 页）}：回到两商品禀赋模型——若 Home 在商品 1 的禀赋上高于世界平均，则 Home 出口商品 1；两国相对价格均等。H-O 中，一国出口密集使用其相对丰裕要素的商品。\redkw{两国的相对要素价格会均等吗？}（红字设问）

\textbf{推导（第 51--52 页，图~\ref{fig:ch07-fpe}）}：贸易使 C 与 S 的相对价格在两国均等。在 PPF 斜率部分有（高亮）：
\[
\frac{P_C}{P_S} = \frac{MPL_S}{MPL_C} = \frac{MPK_S}{MPK_C}
\;\Longrightarrow\;
\left( \frac{w}{r} \right) \text{（手写箭头）}.
\]
因此 \kw{MPL 之比与 MPK 之比在两国均等}。代入 C--D 例子：由 $\frac{P_C}{P_S} = \frac{MPL_S}{MPL_C} = \frac{(1-\beta_S)a_S (K_S/L_S)^{\beta_S}}{(1-\beta_C) a_C (K_C/L_C)^{\beta_C}}$（板书第一步，手写``\kw{相等}''、``1.完全分工→？''批注：注意这里是 both goods produced），再用 $\frac{L_i}{K_i} = \left( \frac{1-\beta_i}{\beta_i} \right) \left( \frac{w}{r} \right)^{-1}$（由 $\frac{w}{r} = \frac{P_i \cdot MPL_i}{P_i \cdot MPK_i}$ 导出，手写标注）得
\[
\frac{P_C}{P_S} = \frac{\beta_S^{\beta_S} (1-\beta_S)^{1-\beta_S} a_S}{\beta_C^{\beta_C} (1-\beta_C)^{1-\beta_C} a_C} \left( \frac{w}{r} \right)^{\beta_S - \beta_C}
\]
（手写：$\beta_S - \beta_C \neq 0 \Rightarrow \frac{w}{r}$ 相等）。结论：\kw{技术相同且两国都生产两种商品时，共同的商品相对价格决定共同的 $w/r$}。

\begin{figure}[H]
  \centering
  \includegraphics[width=0.66\linewidth]{figs/ch07_fig_15.png}
  \caption{第 52 页板书：要素价格均等化的推导——共同的商品相对价格 $\Rightarrow$ 共同的 $w/r$（C--D 情形解析式）。}
  \label{fig:ch07-fpe}
\end{figure}

\begin{intubox}[现实中的要素价格为何不均等？（第 53 页）]
现实中即使两国互相贸易，要素价格通常不均等：不同的 $w/r$（劳动 vs.\ 资本的相对价格）；不同的\kw{技能溢价（skill premium，高亮）}（熟练 vs.\ 非熟练劳动的相对价格）。哪些关键假设在现实中可能不成立？贸易成本为正（相对价格不会完全均等）；国家可能\term{完全专业化}（不再生产两种商品）；各国技术可能不同；商品或要素的质量跨国不同。
\end{intubox}

\begin{intubox}[复习指引（第 54 页）]
回忆中级微观的一般均衡：给定禀赋的交换——Edgeworth 盒。复习：把交换的一般均衡与禀赋模型比较；把生产的一般均衡与 H-O 模型（自给自足）比较。参考：Hal Varian, \textit{Intermediate Microeconomics} (9th ed.), Chapters 32--33。
\end{intubox}

\textbf{H-O 总结（第 55 页）}：即使技术相同，\kw{禀赋差异也能产生贸易}；\redkw{Heckscher--Ohlin 定理}：一国出口密集使用其相对丰裕要素的商品；\redkw{Stolper--Samuelson 定理}：商品相对价格上升提高密集用于其生产的要素的真实回报、降低另一要素的真实回报；附录 2：模型的关键方程结构。

\secn{7.7 附录 1：PPF 斜率的正式推导}

\textbf{关键方程（第 57--59 页）}：从 Stolper--Samuelson 讨论的关键方程出发：
\[
\frac{L_C}{K_C} \cdot \frac{K_C}{\bar{K}} + \frac{L_S}{K_S} \left( 1 - \frac{K_C}{\bar{K}} \right) = \frac{\bar{L}}{\bar{K}} .
\]
想扩张计算机行业（需要收缩鞋业）：$\frac{K_C}{\bar{K}}$ 上升、$\frac{K_S}{\bar{K}}$ 下降。为使左边不变有两种可能：若 $\frac{L_C}{K_C} = \frac{L_S}{K_S}$ 则 LHS 不变——但 H-O 假设 $\frac{L_C}{K_C} < \frac{L_S}{K_S}$（两者都随 $\frac{w}{r}$ 递减）！当计算机行业扩张、$\frac{K_C}{\bar K}$ 上升时，因为 $\frac{L_C}{K_C} < \frac{L_S}{K_S}$，LHS 变小，所以 $\frac{w}{r}$ 必须下降，使 $\frac{L_C}{K_C}$ 与 $\frac{L_S}{K_S}$ 都上升以平衡方程。由于计算机是资本密集的，$K/L$ 下降对其生产的负效应大于鞋（劳动密集）——\kw{扩张计算机生产的机会成本更高，PPF 向下弯（``concave''）}。

\textbf{C--D 数值说明（第 60 页）}：由第 52 页 $\frac{P_C}{P_S} = \frac{\beta_S^{\beta_S}(1-\beta_S)^{1-\beta_S} a_S}{\beta_C^{\beta_C}(1-\beta_C)^{1-\beta_C} a_C} \left( \frac{w}{r} \right)^{\beta_S - \beta_C}$：计算机更资本密集时 $\beta_S - \beta_C < 0$；更多资源（如资本）配置给计算机时，Stolper--Samuelson 推导告诉我们 $\frac{w}{r}$ 下降 $\Rightarrow$ $\left( \frac{w}{r} \right)^{\beta_S - \beta_C}$ 上升 $\Rightarrow$ $\frac{P_C}{P_S}$（PPF 绝对斜率）上升——\kw{向右（横轴为 C）移动时 PPF 变陡！}

\secn{7.8 附录 2：方程结构}

\textbf{生产函数（第 62 页）}：两部门 $i = 1, 2$：$Q_i = A_i L_i^{1-\alpha_i} K_i^{\alpha_i}$；工资与租金率记 $w, r$；设行业 2 资本密集：$0 < \alpha_1 < \alpha_2 < 1$。

\textbf{七组方程（第 63--64 页）}：
\begin{enumerate}
  \item 要素密集度 $\to$ 边际产品：$MPL_i = (1-\alpha_i) A_i \left( \frac{K_i}{L_i} \right)^{\alpha_i}$，$MPK_i = \alpha_i A_i \left( \frac{L_i}{K_i} \right)^{1-\alpha_i}$；
  \item 边际产品 $\to$ 真实回报：$\frac{w}{P_i} = MPL_i$，$\frac{r}{P_i} = MPK_i$；
  \item 生产函数 $\to$ 要素密集度：$\frac{L_i}{K_i} = \left( \frac{1-\alpha_i}{\alpha_i} \right) \left( \frac{w}{r} \right)^{-1}$；
  \item 均衡条件（相对要素供给 $=$ 相对要素需求）：$\frac{\bar{L}}{\bar{K}} = \frac{L_1}{K_1} \frac{K_1}{\bar{K}} + \frac{L_2}{K_2} \frac{K_2}{\bar{K}}$；
  \item 均衡要素价格：$\frac{w}{r} = \frac{\bar{K}}{\bar{L}} \left[ \left( \frac{1-\alpha_1}{\alpha_1} \right) \frac{K_1}{\bar{K}} + \left( \frac{1-\alpha_2}{\alpha_2} \right) \frac{K_2}{\bar{K}} \right]$。
\end{enumerate}

\textbf{方程的机制（第 65--66 页）}：
\begin{enumerate}
  \item 设 Home 相对资本丰裕，则自给自足的 $\frac{P_2}{P_1}$ 低于贸易时；
  \item 开放贸易后 $\frac{P_2}{P_1}$ 上升，$K_2$ 上升、$K_1 = \bar{K} - K_2$ 下降；
  \item $\alpha_2 > \alpha_1 \Rightarrow \frac{1-\alpha_2}{\alpha_2} < \frac{1-\alpha_1}{\alpha_1}$；结合第 2 步与方程 (7) 知 $\frac{w}{r}$ 下降；
  \item $\frac{w}{r}$ 下降时，方程 (5) 意味着 $L_i/K_i$ 对 $i = 1, 2$ 都上升；
  \item $L_i/K_i$ 上升时，方程 (1) 意味着两部门 MPL 下降、方程 (2) 意味着两部门 MPK 上升；
  \item MPL 降、MPK 升（两部门）时，方程 (3) 意味着以两种商品计的真实工资都下降；方程 (4) 意味着以两种商品计的真实资本回报都上升；
  \item 因此：开放贸易后，\kw{相对丰裕要素的所有者明确得益、相对稀缺要素的所有者在真实 terms 上明确受损}。
\end{enumerate}

\secn{7.9 模型构建方法地图（据课后手写笔记整理）}

\begin{mapbox}[SFM 与 H-O 的要素市场对照（课后手写笔记第 67--69 页）]
\textbf{1. 特定要素（SFM）}：$K, T$ 不跨部门配置；$L$ 由市场配置。均衡：$w = P_1 MPL_1 = P_2 MPL_2$（图：$P_1 \cdot MPL_1$ 与 $P_2 \cdot MPL_2$ 两条曲线在 $\bar{L}$ 上相交于 $w^*$，决定 $L_1, L_2$）。贸易的福利影响分析四步（手写）：
\begin{enumerate}
  \item 相对价格变化（假设外生）；
  \item 曲线 $P_1 MPL_1$ vs $P_2 MPL_2$ 的移动（Core：$P_1/P_2 \cdot MPL_1$ vs $MPL_2$）；
  \item $L_1$ 与 $L_2$ 变化；
  \item $\frac{w}{P_1} = MPL_1$ 与 $\frac{w}{P_2} = MPL_2$ 变化。
\end{enumerate}
\textbf{2. Heckscher--Ohlin}：两种要素都可流动。均衡（手写框``Relative labor demand''）：
\[
P_1 MPL_1 = w = P_2 MPL_2, \qquad P_1 MPK_1 = r = P_2 MPK_2
\;\Longrightarrow\;
\frac{MPL_1}{MPL_2} = \frac{w}{r} = \frac{MPL_2}{MPK_2} \text{（两条相对劳动需求曲线）}.
\]
相对供给固定：$\frac{\bar{L}}{\bar{K}} = \frac{L_1}{\bar{K}} + \frac{L_2}{\bar{K}} = \frac{L_1}{K_1} \cdot \frac{K_1}{\bar{K}} + \frac{L_2}{K_2} \cdot \frac{K_2}{\bar{K}}$（手写标注：``share of capital in each sector''；$\frac{L_i}{K_i}$ 为各行业的相对需求曲线）$= D^{\text{avg}}(\frac{w}{r})$——\kw{经济的相对劳动需求}（各行业相对劳动需求按资本份额加权）。图中：固定相对供给（垂直线）与 $D^{\text{avg}}$ 相交定 $w/r$；$D^1, D^2$ 为各行业曲线。
\textbf{3. H-O 中贸易的福利分析六步}（手写）：
\begin{enumerate}
  \item $P_1/P_2$ 变化 $\to$ $\frac{K_1}{\bar{K}}, \frac{K_2}{\bar{K}}$ 变化；
  \item $D^{\text{avg}}$ 曲线\kw{移动}（Shift，高亮）；
  \item $D^{\text{avg}}$ 与 $S$ 相交：$\frac{w}{r}$ 变化；
  \item $\frac{w}{r}$ 变化：$\frac{L_1}{K_1}, \frac{L_2}{K_2}$ 变化（\kw{Move along} $D^1, D^2$ 曲线，高亮）；
  \item $\frac{L_1}{K_1}, \frac{L_2}{K_2}$ 同方向变化 $\Rightarrow$ $MPL_1, MPL_2$ 变化；$MPK_1, MPK_2$ 变化；
  \item $MPL_i = \frac{w}{P_i}$（\kw{工人的福利}）；$MPK_i = \frac{r}{P_i}$（\kw{资本所有者的福利}）。
\end{enumerate}
\end{mapbox}

\begin{figure}[H]
  \centering
  \includegraphics[width=0.5\linewidth]{figs/ch07_fig_16.png}
  \caption{课后手写笔记（一）：SFM 的要素市场——特定要素不流动、劳动由市场配置；$P_1 MPL_1$ 与 $P_2 MPL_2$ 相交定均衡；贸易福利分析四步。}
  \label{fig:ch07-note1}
\end{figure}

\begin{figure}[H]
  \centering
  \includegraphics[width=0.5\linewidth]{figs/ch07_fig_17.png}
  \caption{课后手写笔记（二）：H-O 要素市场——两种要素都可流动；``Relative labor demand''；相对供给固定的分解与 $D^{\text{avg}}$（各行业曲线按资本份额加权）。}
  \label{fig:ch07-note2}
\end{figure}

\begin{figure}[H]
  \centering
  \includegraphics[width=0.5\linewidth]{figs/ch07_fig_18.png}
  \caption{课后手写笔记（三）：H-O 中贸易的福利分析六步——从相对价格变化到工人与资本所有者的福利（Shift vs.\ Move along 的区分）。}
  \label{fig:ch07-note3}
\end{figure}

\begin{sumbox}[本讲总结：通用结论与方法]
\begin{enumerate}
  \item \term{六条假设}：要素跨行业自由流动；两部门要素密集度不同（鞋劳动密集、计算机资本密集，$L_S/K_S > L_C/K_C$ 对一切 $w/r$）；两国相对禀赋不同（$\bar{L}^*/\bar{K}^* > \bar{L}/\bar{K}$）；商品自由贸易、要素不跨国流动；技术相同；同质位似偏好。
  \item \term{三个核心概念辨析}：要素丰裕度是\textbf{相对禀赋比}（不是总量）；要素密集度是\textbf{同一 $w/r$ 下} $L/K$ 的比较（相对要素需求曲线的左右位置，不是使用的总量）；行业内要素比例沿相对要素需求曲线随 $w/r$ 移动（$w/r \uparrow \Rightarrow L/K \downarrow$）。
  \item \term{相对要素需求}：$L_i/K_i = \frac{1-\beta_i}{\beta_i} (\frac{w}{r})^{-1}$（C--D）；经济总体的 $RD = \frac{L_C}{K_C}\frac{K_C}{\bar K} + \frac{L_S}{K_S}\frac{K_S}{\bar K}$ 是各行业劳动密集度的资本份额加权平均，位于两条行业曲线之间；相对供给是垂直线 $\bar{L}/\bar{K}$。
  \item \term{PPF 弯曲与偏斜}：斜率 $= \frac{P_C}{P_S} = \frac{MPL_S}{MPL_C} = \frac{MPK_S}{MPK_C}$；要素密集度不同 $\Rightarrow$ 机会成本递增 $\Rightarrow$ PPF 外凸；禀赋差异 $\Rightarrow$ PPF 向密集使用丰裕要素的商品偏斜。
  \item \term{H-O 定理}：出口密集使用相对丰裕要素的商品；机制与李嘉图同构（自给价格差异 $\to$ 贸易 $\to$ 价格均等化），但原因从技术差异换成禀赋差异（偏好与技术相同假设下）。
  \item \term{Stolper--Samuelson 定理}：$P_C/P_S \uparrow \Rightarrow w/r \downarrow \Rightarrow$ 两行业 $L/K$ 都 $\uparrow \Rightarrow$ 两行业 MPL 都 $\downarrow$、MPK 都 $\uparrow \Rightarrow$ \kw{以两种商品计的真实工资都下降、真实租金都上升}——密集使用涨价商品的要素（K）明确得益、另一要素（L）明确受损，与要素所在行业无关。名义不等式串：$\frac{\Delta w}{w} < \frac{\Delta P_S}{P_S} < \frac{\Delta P_C}{P_C} < \frac{\Delta r}{r}$（Home 情形）。
  \item \term{要素价格均等化}：技术相同 $+$ 两种商品都被生产 $\Rightarrow$ 共同商品价格 $\Rightarrow$ 共同 $w/r$（进而共同的要素绝对价格）；现实中不成立的原因：贸易成本、完全专业化、技术差异、质量差异。
  \item \term{方法论要点}：Shift（相对需求曲线因权重变化而移动）vs.\ Move along（沿行业曲线因 $w/r$ 变化而移动）——手写笔记的六步福利分析是解题模板。
\end{enumerate}
\end{sumbox}
